\item \textbf{El descubrimiento del neutrón por James Chadwick (1932).}
James Chadwick determinó la masa de la partícula recién descubierta disparando neutrones rápidos hacia dos objetivos ligeros de masa conocida ($m_1$ y $m_2$) y midiendo las velocidades de retroceso máximas correspondientes ($v_1$ y $v_2$).
\begin{enumerate}
    \item Asumiendo colisiones elásticas unidimensionales y mecánica clásica, demuestre que la masa del neutrón se obtiene como:
    $$ m_n = \frac{m_1 v_1 - m_2 v_2}{v_2 - v_1} $$
    \item Si al bombardear hidrógeno ($m_1 = 1.00\text{ u}$) se obtuvo $v_1 = 3.3 \times 10^7\text{ m/s}$, y al bombardear nitrógeno ($m_2 = 14.0\text{ u}$) se obtuvo $v_2 = 4.7 \times 10^6\text{ m/s}$, calcule la masa del neutrón calculada por Chadwick en unidades de masa atómica ($\text{u}$).
\end{enumerate}